\section*{Preface}
\addcontentsline{toc}{section}{Preface}

TO DO FOR ME

\begin{enumerate}
    \item Renumber files for Ch. 1-10.  Fix all mains.
    \item Label parts of activities PUE and explain this in opening.  I think keep the whole words for the exercises
    \item Add to this list of notes to the instructor.
    \item (As I go)  create handouts for activities (white = puzzle, then blue... as before.
    \item (As I go) pick new homework and update Moodle
    \item (As I go) relabel quiz question bank and create new quizzes.
    \item (Spring) answers to P, hints to UE.
    \item (Spring) instructor's guides to puzzles (and perhaps some of the activities) including answers?
    \item (Spring) look at Runestone for Levin's book.  Can any of that be of use to me?

\end{enumerate}

To the instructor notes:

\begin{enumerate}
    \item \textbf{class time for activities} A typical activity takes 5-15 minutes. Like the HW, activities are scaffolded: Practice, Understand, and Explore, labeled P, U, and E. Give enough time so that all groups get through the Practice-level problem and, ideally, most groups get through the Understand-level problems. I encourage you to mix up who's working with whom so there aren't groups that always get through everything.  Also can use spies.
    \item \textbf{class time for puzzles}  Allow 15-30 minutes for an activity labeled puzzle. I usually give them 20 minutes and might have some "fast-forward" strategies
    \item \textbf{lecture}  Explain the concept of JBEL (just barely enough lecture).  How do you decide what goes in or not. Backwards design:  what do they need to know in order to do the activity.  The rest they can read examples later for.  To close or not to close?  Consider pointing them to the reading to close, especially after puzzles.  The philosophy is they understand what they do, not what you say.
    \item \textbf{outside of class work}  I have no luck getting my students to do prework, but if you want prework, consider assigning either the first activity, the first part of the first activity, or something like part (a) of all activities.  For homework, I assign a mix of P and U.  I encourage them to do R but don't collect -- we discuss in class the day before the quiz.  I always assign at least one Explore per section and will pick a few more for extra credit.
    \item \textbf{course calendar} in my course that has calculus prereq, we do most sections in a single hour. In our lower level of the course, primarily for information systems and business folks, which only has algebra prereq we cover most sections fully in 2 hours or partially in 1 hour.  
   \item \textbf{section coverage}  Every course has its own goals, so pick and choose. There are 40 sections.  In my class we do all, or nearly all sections by having super short weekly quizzes instead of exams.  For a course with longer quizzes or exams, you'll certainly need to omit some sections.  There are 6 sections marked Excursion that can be skipped without loss of continuity, except the integers mod n is needed for the finite groups subsection. I picked the sections that I typically skip if I run short on time.  Many other sections can be skipped but there might be later homework problems you'll need to avoid. 

   \item \textbf{subsection coverage}  A course that does not have proof writing as one of its learning goals can safely skip all sections named Proof Format. When I am pressed for time I sometimes just assign the activity as part of the later Proof Format subsections as homework and it works fine. A course with a calculus prerequisite can treat the Algebra sections quite lightly.  I usually do not lecture, give them 5 minutes to do the activity and then assign 1-2 exercises as part of the post-work.  In a stronger class one could just assign the exercises or skip altogether.  Any Detour subsections can be skipped without loss of continuity.

\end{enumerate}



What motivated me to create these materials.


\section*{To the student}

What my dream is for you (and a little on how to use the text)

% FROM 105 This textbook is written for you.  Read the narrative examples, listen to examples in class, try the practice exercises (in the accompanying workbook), check with classmates or your instructor, then work the exercises in this textbook, then chat with classmates again, then do more problems, \ldots.  Well, you get the idea:  the best way to learn mathematics is to do it yourself.  I hope you enjoy the course.  I know you will learn a lot of useful algebra.  I believe it will change how you see mathematics.

\section*{To the instructor}

Overview on how to approach this course (and some info on links to instructor tools)

%This textbook is written for students.  That means you won't find list of student learning objectives anywhere, although I'm sure you can infer them from the student-focused ``Do you know \ldots'' questions in each section.  While the narrative examples develop the main theme of each section, I've deliberately left some variations for the practice exercises.   
%
% These practice exercises are designed to be started during class and are printed in a separate workbook for that purpose.   Hand-written solutions to the practice exercises are available (in electronic format) for students to check their work, whether in class or at home.
%
%My greatest success in teaching this course has been to give students room to figure things out for themselves, so try to resist the temptation to show them one of everything.   Listen to your students and help them understand the algebra in their own vocabulary.  You will be impressed.

\section*{Acknowledgments}

Thanks to . . . especially Matt and Jody  also my Matt, Bruce, Tracy, and new gal at Duluth.  Did I get any funding for this work?  Thank MAA Colleagues feedback on many talks, especially the inquiry-based learning community who welcome me in even if I do inquiry differently.  To AIM folks -- converting to Pretext, DOENUT, Runestone, etc.  I was inspired by several Active Texts coming out of GVSU.  Last, my partner Matt Richey for patiently listening to countless details about the course.

%Thanks, first, to the thousands of students who have taken this course.  Their creative approaches to learning mathematics; their unedited criticism and challenge; their often surprising enthusiasm for the course; their patience tried by countless typos and outright mistakes; and their perpetually novel insights have humbled me and challenged everything I thought I knew about teaching and learning mathematics.  They have inspired me time and time again.  I am grateful that they have allowed me to make a difference in their lives.
%
%Thanks, next, to my mathematics colleagues at Augsburg: Mathew Foss (now at North Hennepin Community College), who taught from the very first edition of the textbook back in 1997 and who collaborated in writing earlier versions; Matt Haines, Alyssa Hanson, Rich Flint, and the dozens of other professors who have taught the course over the years from various earlier editions of the textbook; John Zobitz and Jody Sorensen, who edited and created more exercises for this 2012 edition;  and student helpers Ashley Gruhlke and Emma Winegar.
%
%Thanks, also, to my colleagues across campus for allowing us to try something completely different and to my mathematics colleagues nationally for spurring me on.  During the first few years I taught from \emph{Using algebra} by Ethan Bolker.  Much of my approach and probably more examples or exercises than I realize are derived from his vision and from the subsequent text by his colleagues Linda Kime and Judy Clark at University of Massachusetts, Boston.  
%
%A special thanks to Dean Barbara Farley, Augsburg's Center for Teaching and Learning, and Augsburg's Undergraduate Research and Graduate Opportunity program for supporting my work through sabbaticals, travel grants, and summer research grants for both me and students.


